% GENERATED FILE. Edit canonical lessons, then run npm run generate:books.
% canonical-source-hash: fnv1a64:3b9a1b74af5cc627
% canonical-lessons: PA-C224-message-return, PA-C224-message, PA-C224-message-questions, PA-R224-first-pass-saying-no-and-asking, PA-R224-second-pass-past-and-future

\chapter{A Message About Tomorrow}
\label{ch:pa-a2-message}
\hlchaptermodality{224}

\begin{chapteropening}
\textbf{By the end of this chapter:} \emph{I can follow a short message that says what happened, what did not, and what will happen at what time.}

Read a short message and answer by ear what happened, what did not, and when the train will come; then review the whole core.
\end{chapteropening}

\section[\texorpdfstring{\pa{ਸੁਨੇਹਾ} · \pa{ਨੌਂ} \pa{ਵਜੇ} (sunehā · nauṁ vaje)}{ਸੁਨੇਹਾ · ਨੌਂ ਵਜੇ (sunehā · nauṁ vaje)}]{Return to a message and a time}

\label{lesson:PA-C224-message-return}



\begin{quote}
\emph{Hear:} \emph{sunehā} Say \textbf{\pa{ਸੁਨੇਹਾ}}, \textquotedblleft{}a message\textquotedblright{}. Then \emph{at nine}: \textbf{\pa{ਨੌਂ} \pa{ਵਜੇ}}.
\end{quote}

\begin{grammarlens}[title={Grammar Lens}]
A message from a friend often says three things: what happened, what did not, and what will happen when. You can say all three now. Listen for \textbf{\pa{ਕੱਲ੍ਹ}}: the verb tells you whether it is yesterday or tomorrow.
\end{grammarlens}

\subsection*{Your turn}
\begin{itemize}
\raggedright
  \item \emph{Hear:} \emph{a message} Say \textbf{\pa{ਸੁਨੇਹਾ}}.
  \item \emph{Hear:} \emph{at eight} Say \textbf{\pa{ਅੱਠ} \pa{ਵਜੇ}}.
  \item \emph{Hear:} \emph{I went to the market yesterday} Say it with the ending that fits you.
  \item \emph{Hear:} \emph{we will go tomorrow} Say \textbf{\pa{ਅਸੀਂ} \pa{ਕੱਲ੍ਹ} \pa{ਜਾਵਾਂਗੇ}}.
\end{itemize}

\begin{tcolorbox}[breakable,colback=teal!4,colframe=teal!35!black,title={Before you move on}]
\textbf{\pa{ਕੱਲ੍ਹ}} with \textbf{\pa{ਗਿਆ}}? (Yesterday.) With \textbf{\pa{ਜਾਵਾਂਗੇ}}? (Tomorrow.)
\end{tcolorbox}
\section[(reading)]{A message about tomorrow}

\label{lesson:PA-C224-message}



\begin{quote}
Say \emph{a message}. Then \emph{Aman}.

Aman sends you a message. Read it for one thing first: what time you have to be at the station.
\end{quote}

\begin{tcolorbox}[breakable,skin=enhanced,colback=orange!6,colframe=orange!45!black,boxrule=0.5pt,arc=1mm,left=6pt,right=6pt,top=4pt,bottom=4pt,fonttitle=\bfseries,title={Read this}]
\begin{quote}
\textbf{\pa{ਸਤਿ} \pa{ਸ੍ਰੀ} \pa{ਅਕਾਲ} \pa{ਜੀ}!}
\end{quote}

\begin{quote}
\textbf{\pa{ਮੈਂ} \pa{ਅਮਨ} \pa{ਹਾਂ।}} \textbf{\pa{ਕੱਲ੍ਹ} \pa{ਮੈਂ} \pa{ਬਾਜ਼ਾਰ} \pa{ਗਿਆ।}} \textbf{\pa{ਮੈਂ} \pa{ਇੱਕ} \pa{ਕਿਤਾਬ} \pa{ਖ਼ਰੀਦੀ}, \pa{ਪਰ} \pa{ਮੈਂ} \pa{ਚਾਹ} \pa{ਨਹੀਂ} \pa{ਪੀਤੀ।}} \textbf{\pa{ਕੱਲ੍ਹ} \pa{ਸਵੇਰੇ} \pa{ਅਸੀਂ} \pa{ਰੇਲ} \pa{ਨਾਲ} \pa{ਜਾਵਾਂਗੇ।}} \textbf{\pa{ਰੇਲ} \pa{ਨੌਂ} \pa{ਵਜੇ} \pa{ਆਵੇਗੀ।}} \textbf{\pa{ਮੈਂ} \pa{ਅੱਠ} \pa{ਵਜੇ} \pa{ਸਟੇਸ਼ਨ} \pa{ਜਾਵਾਂਗਾ।}} \textbf{\pa{ਟਿਕਟ} \pa{ਨਾ} \pa{ਭੁੱਲੋ}!} \textbf{\pa{ਸ਼ਾਇਦ} \pa{ਮੇਰਾ} \pa{ਭਰਾ} \pa{ਆਵੇਗਾ।}} \textbf{\pa{ਫਿਰ} \pa{ਮਿਲਾਂਗੇ।}}
\end{quote}

\emph{Read it:} the message once, straight through, without translating

\emph{Read it:} it a second time, and find every line with a clock time
\end{tcolorbox}

\subsection*{What to know first}
\textbf{\pa{ਕੱਲ੍ਹ}} comes twice: with \textbf{\pa{ਗਿਆ}} it is yesterday, with \textbf{\pa{ਜਾਵਾਂਗੇ}} tomorrow. \textbf{\pa{ਟਿਕਟ} \pa{ਨਾ} \pa{ਭੁੱਲੋ}!} is a request built like \textbf{\pa{ਨਾ} \pa{ਉੱਠੋ}}, on \textbf{\pa{ਭੁੱਲਣਾ}}, \textquotedblleft{}to forget\textquotedblright{}. The last line but one is the only one that is not sure.

\subsection*{Your turn}
\begin{itemize}
\raggedright
  \item \emph{Answer:} what did Aman buy yesterday? (A book: \textbf{\pa{ਮੈਂ} \pa{ਇੱਕ} \pa{ਕਿਤਾਬ} \pa{ਖ਼ਰੀਦੀ।}})
  \item \emph{Answer:} what did he not do? (He didn't drink tea.)
  \item \emph{Answer:} when will the train come? (At nine: \textbf{\pa{ਨੌਂ} \pa{ਵਜੇ}}.)
  \item \emph{Answer:} who will perhaps come? (His brother.)
\end{itemize}

\begin{tcolorbox}[breakable,colback=teal!4,colframe=teal!35!black,title={Before you move on}]
When will Aman be at the station? (At eight.) What must you not forget? (The ticket.)
\end{tcolorbox}
\section[\texorpdfstring{\pa{ਰੇਲ} \pa{ਕਦੋਂ} \pa{ਆਵੇਗੀ}? (rel kadoṁ āvegī?)}{ਰੇਲ ਕਦੋਂ ਆਵੇਗੀ? (rel kadoṁ āvegī?)}]{The message, by ear}

\label{lesson:PA-C224-message-questions}



\begin{quote}
Hold Aman's message in mind: yesterday, and tomorrow.
\end{quote}

\subsection*{Your turn}
\begin{itemize}
\raggedright
  \item \emph{Hear:} \emph{Aman kall kitthe giā?} Answer: \textbf{\pa{ਉਹ} \pa{ਬਾਜ਼ਾਰ} \pa{ਗਿਆ।}}
  \item \emph{Hear:} \emph{kī Aman ne chāh pītī?} Answer: \textbf{\pa{ਨਹੀਂ} \pa{ਜੀ}, \pa{ਉਸ} \pa{ਨੇ} \pa{ਚਾਹ} \pa{ਨਹੀਂ} \pa{ਪੀਤੀ।}}
  \item \emph{Hear:} \emph{rel kadoṁ āvegī?} Answer: \textbf{\pa{ਰੇਲ} \pa{ਨੌਂ} \pa{ਵਜੇ} \pa{ਆਵੇਗੀ।}}
  \item \emph{Hear:} \emph{Aman kinne vaje saṭeshan jāvegā?} Answer: \textbf{\pa{ਅੱਠ} \pa{ਵਜੇ।}}
\end{itemize}

\begin{tcolorbox}[breakable,colback=teal!4,colframe=teal!35!black,title={Before you move on}]
Tell the message back in three lines: what Aman did, what he didn't, and what will happen tomorrow.
\end{tcolorbox}
\section[\texorpdfstring{\pa{ਨਹੀਂ} · \pa{ਨਾ} · \pa{ਕੀ} (nahīṁ · nā · kī)}{ਨਹੀਂ · ਨਾ · ਕੀ (nahīṁ · nā · kī)}]{Looking back: saying no and asking}

\label{lesson:PA-R224-first-pass-saying-no-and-asking}



\begin{quote}
Say \emph{I don't drink tea}. Then ask \emph{do you drink tea?}
\end{quote}

\subsection*{Your turn}
\begin{itemize}
\raggedright
  \item \emph{Hear:} \emph{I go to school every day} Say it with the ending that fits you.
  \item \emph{Hear:} \emph{we eat roti every day} Say \textbf{\pa{ਅਸੀਂ} \pa{ਰੋਜ਼} \pa{ਰੋਟੀ} \pa{ਖਾਂਦੇ} \pa{ਹਾਂ।}}
  \item \emph{Hear:} \emph{the garden is not open} Say \textbf{\pa{ਬਾਗ਼} \pa{ਖੁੱਲ੍ਹਾ} \pa{ਨਹੀਂ} \pa{ਹੈ।}}
  \item \emph{Hear:} \emph{please sit here; don't get up yet} Say \textbf{\pa{ਇੱਥੇ} \pa{ਬੈਠੋ।} \pa{ਅਜੇ} \pa{ਨਾ} \pa{ਉੱਠੋ।}}
  \item \emph{Hear:} \emph{are you a traveller? No.} Say \textbf{\pa{ਕੀ} \pa{ਤੁਸੀਂ} \pa{ਯਾਤਰੀ} \pa{ਹੋ}? \pa{ਨਹੀਂ} \pa{ਜੀ}, \pa{ਮੈਂ} \pa{ਯਾਤਰੀ} \pa{ਨਹੀਂ} \pa{ਹਾਂ।}}
  \item \emph{Read it:} \textbf{\pa{ਕੀ} \pa{ਉਹ} \pa{ਰੋਜ਼} \pa{ਸਕੂਲ} \pa{ਜਾਂਦੀ} \pa{ਹੈ}?} Give its meaning.
\end{itemize}

\begin{tcolorbox}[breakable,colback=teal!4,colframe=teal!35!black,title={Before you move on}]
\textbf{\pa{ਨਹੀਂ}} tells, \textbf{\pa{ਨਾ}} asks someone not to, \textbf{\pa{ਕੀ}} at the front asks yes or no. Say one sentence for each.
\end{tcolorbox}
\section[\texorpdfstring{\pa{ਗਿਆ} · \pa{ਪੀਤੀ} · \pa{ਜਾਵਾਂਗੇ} (giā · pītī …)}{ਗਿਆ · ਪੀਤੀ · ਜਾਵਾਂਗੇ (giā · pītī …)}]{Looking back: yesterday and tomorrow}

\label{lesson:PA-R224-second-pass-past-and-future}



\begin{quote}
Say \emph{Aman went home yesterday}. Then \emph{we will go tomorrow}.
\end{quote}

\subsection*{Your turn}
\begin{itemize}
\raggedright
  \item \emph{Hear:} \emph{a woman came yesterday} Say \textbf{\pa{ਔਰਤ} \pa{ਕੱਲ੍ਹ} \pa{ਆਈ।}}
  \item \emph{Hear:} \emph{I got up early and drank tea} Say it with the ending that fits you.
  \item \emph{Hear:} \emph{Aman bought a book and she gave money} Say \textbf{\pa{ਅਮਨ} \pa{ਨੇ} \pa{ਕਿਤਾਬ} \pa{ਖ਼ਰੀਦੀ} \pa{ਅਤੇ} \pa{ਉਸ} \pa{ਨੇ} \pa{ਪੈਸਾ} \pa{ਦਿੱਤਾ।}}
  \item \emph{Hear:} \emph{I took out the book} Say \textbf{\pa{ਮੈਂ} \pa{ਕਿਤਾਬ} \pa{ਕੱਢੀ।}}
  \item \emph{Hear:} \emph{what will you do tomorrow?} Say \textbf{\pa{ਤੁਸੀਂ} \pa{ਕੱਲ੍ਹ} \pa{ਕੀ} \pa{ਕਰੋਗੇ}?}
  \item \emph{Hear:} \emph{perhaps we will go next week} Say \textbf{\pa{ਸ਼ਾਇਦ} \pa{ਅਸੀਂ} \pa{ਅਗਲੇ} \pa{ਹਫ਼ਤੇ} \pa{ਜਾਵਾਂਗੇ।}}
  \item \emph{Read it:} \textbf{\pa{ਰੇਲ} \pa{ਨੌਂ} \pa{ਵਜੇ} \pa{ਆਵੇਗੀ।}} Give its meaning, then say when Aman will be at the station.
  \item \emph{Write it:} cover the models and write \textbf{\pa{ਮੈਂ} \pa{ਕੱਲ੍ਹ} \pa{ਸਕੂਲ} \pa{ਨਹੀਂ} \pa{ਗਿਆ।}} or its woman's form
\end{itemize}

\begin{tcolorbox}[breakable,colback=teal!4,colframe=teal!35!black,title={Before you move on}]
Tell your yesterday in two lines and your tomorrow in two lines.
\end{tcolorbox}
